\section{Simply-typed \WLOG: further details}
\applabel{simply-typed wlog details}

In this Appendix, we provide further details to \secref{wlog construction}.
Our first proof obligation is a comment we immediately before \exampleref{tree example}:

\begin{lemma}[well-known]\lemmalabel{action exponential}
  Let $\A$ and $\B$ be $\G$-actions. The exponential $\A \to \B$ in $\Action\G$ is given by the set of functions between their carriers:
  \begin{gather*}
  \carrier{\A \to \B}\definedby \carrier\A \to \carrier\B
  \qquad
  (g \act[\A \to \B] f) \definedby \fun a g \act[\B] f(\inv g\act[\A]a)
  \\
  \eval : (\A \to \B) \times \A \equito \B
  \qquad
  \eval \pair fa \definedby f\,a
  \end{gather*}
\end{lemma}
\begin{proof}
  We indeed have an action. Unit preservation:
  \[
  (\unit[\G] \act[\A \to \B] f)a
  = \unit[\G]\act[\B] f(\inv{\unit[\G]}\act a) = f\,a
  \implies
  \unit[\G] \act[\A \to \B] f = f
  \]
  Action preservation:
  \begin{align*}
  (g \act[\A \to \B] h \act[\A\to\B] f)a
  &=
  g \act[\B] (h \act[\A\to\B] f)(\inv g \act[\A] a)
  =
  g \act[\B] (h \act[\B] f(\inv h \act[\A] (\inv g \act[\A] a)))
  \\&=
  (g \mul[\G] h) \act[\B] f\parent{\inv*{\smash{g \mul[\G] h}}\act[\A] a}
  =
  \parent{(g \mul[\G] h) \act[\A \to \B] f} a
  \\&\implies
  (g \act[\A \to \B] h \act[\A\to\B] f) = (g \mul[\G] h) \act[\A \to \B] f
  \end{align*}

  Evaluation is equivariant. Calculate:
  \begin{align*}
  \eval(g \act[(\A\to\B)\times \A] \pair fa)
  &=
  \eval\pair{g\act[\A\to\B] f}{g\act[\A] a}
  =
  (g\act[\A\to\B]f)\,(g\act[\A] a)
  \\&=
  g \act[\B] (f(\inv g \act[A] (g \act[\A] a)))
  =
  g \act[\B] (f\,a)
  =
  g \act[\B] \eval\pair fa
  \tag*{\qedhere}
  \end{align*}

  Given any equivariant map $h : \C \times \A \equito \B$, its curried
  map is also equivariant:
  \begin{align*}
    (g \act[\A \To \B] \curry h c)a
    &=
    g \act[\B] \curry h (c; \inv g \act[\A] a)
    =
    g \act[\B] h (c, \inv g \act[\A] a)
    \explain={equivariance}
    h (g \act[\C] c, g \act \inv g \act[\A] a)
    \\&=
    h (g \act[\C] c, a)
    =
    (\curry h (g \act[\C]c) a
    \\&\implies\quad
    g \act[\A \to \B] \curry h c = \curry h (g \act[\C] c)
  \end{align*}
  and so $\curry h : \C \equito (\A \to \B)$. The curried function
  lifts $h$ along evaluation. Now if we have an equivariant map $g :
  \C \equito (\A \to \B)$ that lifts $h$ along evaluation, then by the
  set-theoretic argument, it coincides with $\curry h$.
\end{proof}

Next, we characterise the canonical elements of
canonisers of the form $\carrier \A \times \carrier \B \xto{\projection_1}\A \xto\canoniser \carrier\G$:

\begin{proof}[of \lemmaref{projection canoniser}]%
  Take an arbitrary canonical element $(\canoniser \compose \projection_1) \seq{a, b}
  \act \seq{a, b} \in \Canon{\parent{\canoniser \compose
  \projection_1}}$. We then have
  \[
    (\canoniser \compose \projection_1) \seq{a, b} \act \seq{a, b}
    = \canoniser a \act \seq{a, b}
    = \seq{\canoniser a \act a, \canoniser a \act b}
    \in \Canon\canoniser \times \carrier\B
  \]
  Conversely, given an arbitrary pair $\seq{\canoniser a \act a, b} \in
  \Canon\canoniser \times \carrier\B$, let $b' = \inv*{\canoniser a} \act b$.
  We then have:
  \[
    \seq{\canoniser a \act a, b}
    = \seq{\canoniser a \act a, (\canoniser a \mul \inv*{\canoniser a}) \act b}
    = \canoniser a \act \seq{a, b'}
    = ((\canoniser \compose \projection_1) \seq{a, b'}) \act \seq{a, b'}
    \in \Canon{\parent{\canoniser \compose \projection_1}}
  \]
  Therefore $\Canon{\parent{\canoniser \compose
  \projection_1}} = \Canon\canoniser \times \carrier\B$.
\end{proof}

Next, we proof the main equivariance theorem for the \WLOG{} construction:
\begin{proof}[of \theoremref{wlog-construction}]%
  Let us first show that $\wlg(c, f)$ is equivariant. Take arbitrary
  $g \in \carrier\G$ and $a \in \carrier\A$. By definition and by the
  group and action axioms, we can rewrite:
  \begin{align*}
    \wlg(c, f)(g \act a)
    &=
    \inv{c (g \act a)} \act f(c (g \act a) \act (g \act a)) \\
    &=
    \inv{c (g \act a)} \act f(c (g \act a) \act (g \act \inv{(c a)} \act c a \act a)) \\
    &=
    \inv{c (g \act a)} \act f\big((c (g \act a) \mul g \mul \inv{(c a)}) \act (c a \act a)\big) \\
    \intertext{
      Since both $c a \act a$ and $c (g \act a) \act (g \act a)$ belong to $\Canon{\canoniser}$, equivariance of $f$ allows us to rewrite:
    }
    &=
    \inv{c (g \act a)} \act \big((c (g \act a) \mul g \mul \inv{(c a)}\big) \act f(c a \act a)) \\
    &=
    g \act (\inv{(c a)} \act f(c a \act a)) \\
    &=
    g \act \wlg(c, f)\,a
  \end{align*}
  as desired.

  Next, let us show that $\wlg(c, f)$ extends $f$. Take an arbitrary
  canonical element $a \in \Canon\canoniser$. Since $c a \act a$
  also lies in $\in \Canon\canoniser$, the action axioms imply:
  \[
    \wlg(c, f)\,a
    =
    \inv{(c a)} \act f(c a \act a)
    \explain={$f$ is equivariant}
    \inv{(c a)} \act c a \act f a
    =
    f a
  \]

  Finally, take an arbitrary $\phi : \A \equito \B$ that extends $f$
  along $\Canon\canoniser\subset\carrier\A$. Then, for an arbitrary
  $a \in \carrier\A$, we have
  \[
    \phi a
    =
    \phi (\inv{(c a)} \act c a \act a)
    \explain={$\phi$ equivariant}
    \inv{(c a)} \act \phi (c a \act a)
    \explain={$\phi$ extends $f$}
    \inv{(c a)} \act f (c a \act a)
    \explain={by definition}
    \wlg(c, f)\,a
  \]
  as desired.
\end{proof}

The order of the next few results
differs from the order in which they appear in the main text.
They concern the characterisation of
strictly equivariant and immediate canonisers from
\subsecref{wlog-construction}.
Strict equivariance seem to be a more natural concept to the
audience when we present this work, so the main text starts there.
Here we first characterise immediate canonisers, and use this
characterisation as part of the proof. This order make it easier to
see that we avoid circular reasoning.

We start with immediate canonisers as every strictly equivariant
canoniser is also immediate:
\begin{proof}[of \propositionref{strictly equivariant canonisers are immediate}]%
  Calculate:
  \[
  \canoniser(\canoniser a \mul a)
  \explain={equivariance}
  \canoniser a \mul \inv*{\canoniser a}
  =
  \unit
  \]
  as desired.
\end{proof}

We characterise immediate canonisers
via the condition $\Canon\canoniser =
\inv\canoniser[\unit]$:

\begin{proof}[of \propositionref{immediate canonical is kernel}]%
  Take an arbitrary $\canoniser a \act a \in \Canon\canoniser$. Since
  $\canoniser$ is immediate, we have $\canoniser \parent{\canoniser a \act a}
  = \unit$, and so $\canoniser a \act a \in \inv{\canoniser}[\unit]$,
  thus $\Canon\canoniser \subseteq \inv{\canoniser}[\unit]$. The
  converse inclusion always holds. Indeed, if $a \in \inv{\canoniser}[\unit]$,
  then $\canoniser a = \unit$ and so $a = \canoniser a \act a \in
  \Canon\canoniser$.

  Conversely, let us show that $\canoniser$ is immediate and take an
  arbitrary $a \in \carrier\A$. We have $\canoniser a \act a \in
  \Canon\canoniser$ by definition, thus $\canoniser a \act a \in
  \inv{\canoniser}[\unit]$ by assumption, and so
  $\canoniser(\canoniser a \act a) = \unit$ as required.
\end{proof}



The next result characterises strict equivariance of a canoniser via the
equivariance of every core function. The proof uses the concept of
stabilisers that we introduce in \subsecref{equivariant canonisers}.
The proof is dependently-typed in nature, and it is best understood in
the context of \secref{dependent types} and \subsecref{equivariant
  canonisers} in particular. However, we make the proof self-contained
and re-introduce the concepts that are relevant for it.

\begin{proof}[of \theoremref{all core functions are partially equivariant}]%
  First, take a strictly equivariant canoniser $\canoniser :
  \carrier\A \to \carrier\G$, and let $\B$ be an arbitrary action. We
  need to show that every core function $f : \Canon\canoniser \to
  \carrier\B$ is equivariant. Take $a, g\act a \in
  \Canon\canoniser$. Since $\canoniser$ is strictly equivariant,
  by \propositionref{strictly equivariant canonisers are
  immediate}, it is an immediate canoniser. So by \propositionref{immediate canonical is kernel}, $\canoniser a = \unit =
  \canoniser (g \act a)$, therefore:
  \[
    \inv g = \canoniser a \mul \inv g
    \explain*={$\canoniser$ strictly equivariant}
    \canoniser\parent{g \act a}
    = \unit
  \]
  and equivariance holds vacuously.

  Conversely, assume that every core function is equivariant, and let us show that $\canoniser$ is strictly equivariant.
  Recall that for every $a \in \carrier\A$, we define the stabiliser sub-group of $a$ as the set of symmetries that fix $a$:
  \[
    \stabilise a \definedby \set{s \in \carrier\G \suchthat s \act a = a}
  \]
  Take the following action over the
  the set of dependent pairs consisting of an
  element in $\carrier\A$ and one
  of its stabilising symmetries:
  \[
    \carrier\B \definedby (a : \carrier\A) \times \stabilise a
    \qquad
    g \act \pair as \definedby \pair{g \act a}{s \mul \inv g}
  \]
  Indeed, this definition is well-typed because:
  \[
    (s \mul \inv g)\mul (g \act a)
    =
    s \act (\inv g \mul g) \act a
    =
    s \act a
    =
    a
  \]
  It is also an action, since:
  \begin{gather*}
    \unit \act \pair as =
    \pair{\unit \act a}{s \act \inv\unit} =
    \pair as
    \\
    g \act h \act \pair as
    =
    \pair{g \act h \act a}{s \mul \inv h \mul \inv g}
    =
    \pair{(g \mul h) \act a}{s \mul \inv*{g \mul h }}
    =
    (g\mul h)\act \pair as
  \end{gather*}
  Define:
  \[
    f : \Canon\canoniser \to \carrier\B
    \qquad
    f a \definedby \pair a\unit
  \]
  By our assumption, its \WLOG extension is equivariant: $\wlg(c,f) :
  \A \equito \B$. Note that this map is given by:
  \[
    \wlg(f,c)a = \inv*{\canoniser a} \act f\parent{\canoniser a\mul a}
    = \inv*{\canoniser a} \act \pair{\canoniser a\mul a}\unit
    = \pair{a}{\inv*{\inv*{\canoniser a}}}
    = \pair{a}{\canoniser a}
  \]
  Now take any $a \in \carrier\A$ and $g \in
  \carrier\G$, and calculate:
  \begin{align*}
    \pair{g\act a}{\canoniser\parent{g\act a}}
    =
    \wlg(c,f)(g\act a)
    =
    g\act \wlg(c,f)a
    =
    g\act \pair{a}{\canoniser a}
    =
    \pair{g \act a}{\parent{\canoniser a}\mul \inv g}
  \end{align*}
  Comparing the two second components, we have
  \(
    \canoniser\parent{g \act a} = (\canoniser a)\mul \inv g
  \) and the canoniser is strictly equivariant.
\end{proof}

The next result shows that actions that admit strictly equivariant
canonisers must be semi-regular, which makes this condition too strong
for some actions:

\begin{proof}[of \lemmaref{semi-regularity}]%
  Assume $g \act a = a$. By equivariance of the canoniser, we thus have
  $\canoniser a = \canoniser\parent{g \act a} = \canoniser a \mul \inv g$ and by cancelling
  $\canoniser a$ on the left, we get $\unit = g$.
\end{proof}

The final result concerns inheriting strict equivariance and
immediateness from a canoniser:

\begin{proof}[of \lemmaref{equivariance of projection canonisers}]%
  To show strict equivariance, we have that the projection is an equivariant map $\projection_1 : \A \times \B
  \equito \A$, and equivariant maps compose.
%
  For immediateness, we calculate
    \begin{align*}
    \canoniser\compose\projection_1\parent{(\canoniser\compose\projection_1{\pair ab})\mul{\pair ab}}
    &=
    \canoniser\compose\projection_1\parent{\canoniser a\mul{\pair ab}}
    =
    \canoniser\compose\projection_1\parent{\pair {\canoniser a\mul a}{\canoniser a\mul b}}
    \\&=
    \canoniser\parent{\canoniser{a} \mul a}
    =
    \unit
    \end{align*}
    and so $\canoniser \compose \projection_1$ is immediate when $\canoniser$ is.
\end{proof}
